\section{Research continuum balances and MLS-MPM discretization}
\label{sec:mpm}
Research fixtures solve \emph{one} continuum phase (paste or liquid), optionally
coupled to one rigid pellet. They do not yet solve a mixture of water, fines and pellets.
For current density $\rho$, velocity $\mathbf{v}$ and Cauchy stress
$\boldsymbol\sigma$, the local balances motivating the discretization are
\begin{equation}
 \frac{D\rho}{Dt}+\rho\nabla\cdot\mathbf{v}=0,\qquad
 \rho\frac{D\mathbf{v}}{Dt}=\nabla\cdot\boldsymbol\sigma+\rho\mathbf{g}.
 \label{eq:continuum-balance}
\end{equation}
The reference-volume stress is Kirchhoff stress,
$\boldsymbol\tau=J\boldsymbol\sigma$, where $J$ is current volume divided by
reference volume. The constitutive history is phase dependent: paste stores an
elastic deformation tensor; liquid stores a scalar volume ratio. Neither
phase stores the full multiphase state of the intended porous material.

\subsection{Particle and grid variables}
A material point $p$ stores $m_p$, $V_{0p}$, $\mathbf{x}_p$, $\mathbf{v}_p$,
an affine velocity matrix $\mathbf{C}_p$, stress and constitutive history.
The regular grid has spacing $h$ and positions $\mathbf{x}_i$. Material points
carry history across steps; grid mass and momentum are rebuilt each step.
The affine representation approximates local velocity as
\begin{equation}
 \mathbf{v}(\mathbf{x})\simeq
 \mathbf{v}_p+\mathbf{C}_p(\mathbf{x}-\mathbf{x}_p).
\end{equation}
This reduces the loss of rotational/local linear motion characteristic of plain
particle-in-cell transfer \cite{apic}. It does not guarantee energy conservation.

Let $q=|r|$ for normalized distance $r=(x_i-x_p)/h$. The quadratic B-spline is
\begin{equation}
 N(r)=\begin{cases}
  \tfrac34-r^2,&0\leq q<\tfrac12,\\
  \tfrac12(\tfrac32-q)^2,&\tfrac12\leq q<\tfrac32,\\
  0,&q\geq\tfrac32.
 \end{cases}\qquad
 w_{ip}=\prod_{\alpha=1}^3N\!\left(\frac{x_{i\alpha}-x_{p\alpha}}{h}\right).
 \label{eq:weights}
\end{equation}
Each point uses a $3^3=27$ node stencil. With complete support,
\begin{equation}
 \sum_iw_{ip}=1,\quad \sum_iw_{ip}\mathbf{d}_{ip}=\mathbf{0},\quad
 \sum_iw_{ip}\mathbf{d}_{ip}\mathbf{d}_{ip}^{T}=\tfrac{h^2}{4}\mathbf{I},
 \qquad\mathbf{d}_{ip}=\mathbf{x}_i-\mathbf{x}_p. \label{eq:moment-identities}
\end{equation}
The last identity supplies the $4/h^2$ MLS factor, rather than a separately
coded derivative of each shape function. Boundary padding preserves a full
stencil for valid particles; leaving that padded grid is not silently accepted
as an open boundary in the current closed wire fixtures.

\begin{figure}[htbp]
\centering
\begin{tikzpicture}
\begin{axis}[width=.49\textwidth,height=5.1cm,xlabel={$r=(x_i-x_p)/h$},
 ylabel={$N(r)$},xmin=-1.6,xmax=1.6,ymin=0,ymax=.85,
 xtick={-1.5,-.5,.5,1.5},grid=major,domain=-1.5:1.5,samples=151]
 \addplot[blue!75!black,thick] {abs(x)<0.5 ? 0.75-x*x : 0.5*(1.5-abs(x))^2};
\end{axis}
\begin{scope}[shift={(10,1.8)},scale=.85,>=Latex,font=\small]
 \draw[step=1,gray!45] (-1.3,-1.3) grid (1.4,1.4);
 \foreach \x in {-1,0,1}{\foreach \y in {-1,0,1}{\fill[blue!60] (\x,\y) circle (.065);}}
 \fill[orange!90!black] (.23,.35) circle (.09);
 \foreach \x in {-1,0,1}{\foreach \y in {-1,0,1}{\draw[orange!60,thin] (.23,.35)--(\x,\y);}}
 \node[anchor=west] at (.45,.55) {$p$};
 \draw[<->] (-1,-1.65)--node[below]{$h$}(0,-1.65);
 \node[align=center] at (0,2) {2D stencil slice\\9 nodes shown; 27 in 3D};
\end{scope}
\end{tikzpicture}
\caption{Analytical interpolation kernel and an illustrative stencil slice.
These are method diagrams, not simulation output. Exactly zero-weight stencil
entries must not create duplicate active grid nodes.}
\label{fig:kernel}
\end{figure}

\subsection{Particle-to-grid transfer and explicit stress force}
The implemented MLS-MPM transfer \cite{mlsmpm} combines affine momentum and an
explicit stress impulse:
\begin{align}
 m_i&=\sum_pw_{ip}m_p,\\
 \mathbf{p}_i^*&=\sum_pw_{ip}\left[m_p\mathbf{v}_p+
 \left(m_p\mathbf{C}_p-\Delta t\frac4{h^2}V_{0p}\boldsymbol\tau_p\right)
          \mathbf{d}_{ip}\right], \label{eq:p2g}\\
 \mathbf{v}_i^*&=\frac{\mathbf{p}_i^*}{m_i}+\Delta t\mathbf{g},\quad m_i>0.
\end{align}
In weak-form language the stress contribution is
$-\Delta t\sum_pV_{0p}\boldsymbol\tau_p\nabla w_{ip}$ with the MLS gradient
surrogate $\nabla w_{ip}\simeq(4/h^2)w_{ip}\mathbf{d}_{ip}$.
Dimensional check: $V_0\tau$ has units of energy; multiplying by
$\Delta t\,\mathbf{d}/h^2$ gives momentum.
Equation~\eqref{eq:moment-identities} cancels internal stress impulses in the
sum over nodes. That cancellation can fail in bookkeeping if an active node
is processed twice; the zero-deposit regression explicitly guards this case.

\subsection{Wall treatment, transfer back, and local update}
At a wall with inward unit normal $\mathbf{n}$, let
$v_n=\mathbf{v}^*\cdot\mathbf{n}$ and
$\mathbf{t}=\mathbf{v}^*-v_n\mathbf{n}$. Only $v_n<0$ is constrained:
\begin{equation}
 \mathbf{v}^{N}=\mathbf{t},\qquad
 \mathbf{v}^{+}=\mathbf{t}
 \max\left(0,1-\frac{\mu(-v_n)}{\|\mathbf{t}\|}\right), \label{eq:wall-projection}
\end{equation}
with the sticking/zero-tangent case handled explicitly. Axis-aligned wall
constraints are applied sequentially at corners. This is a nodal velocity
projection with Coulomb tangential reduction. It is \emph{not} an incompressible
pressure-Poisson solve and does not enforce a pointwise watertight surface.

After wall and optional pellet coupling updates,
\begin{align}
 \mathbf{v}_p^{n+1}&=\sum_iw_{ip}\mathbf{v}_i^+,\\
 \mathbf{C}_p^{n+1}&=\frac4{h^2}\sum_iw_{ip}\mathbf{v}_i^+
                           \mathbf{d}_{ip}^{T}, \label{eq:g2p}\\
 \mathbf{x}_p^{n+1}&=\mathbf{x}_p^n+\Delta t\,\mathbf{v}_p^{n+1}.
\end{align}
Paste uses the trial elastic tensor
$\mathbf{F}_e^{\rm tr}=(\mathbf{I}+\Delta t\mathbf{C}_p^{n+1})\mathbf{F}_e^n$;
liquid advances its scalar $J$ with $\operatorname{tr}\mathbf{C}_p^{n+1}$.
The constitutive operations are given in Section~\ref{sec:constitutive}.
A locally implicit plastic return does not make the global momentum step
implicit: the stress in \eqref{eq:p2g} is the stored stress from the preceding
accepted state.

\subsection{Adaptive integration and transactional state}
The research solver forms a heuristic upper bound
\begin{equation}
 \Delta t_{\rm stab}=\min\left\{
 \frac{\mathrm{CFL}\,h}{c+\max_p\|\mathbf{v}_p\|},\quad
 \frac{h^2}{4\nu},\quad 0.2\sqrt{m_r/k_n},\quad
 \frac{\mathrm{CFL}\,h}{\|\mathbf{v}_r\|}\right\}, \label{eq:mpm-step}
\end{equation}
where viscous and rigid terms apply only when present/nonzero. Here $c$ is the
material wave speed and $\nu$ the implementation's kinematic diffusivity.
A recovery factor can reduce the candidate step further. Endpoint clipping
then lands on the requested record/final time. Nonfinite state, constitutive
failure or excessive displacement rejects the trial; the solver restores
physical history and ledgers, halves the recovery scale and retries. A minimum
step failure is reported instead of silently changing the physics.

The checkpoint persists limiter/rejection history. A timestep-refinement case
is not accepted merely because one step happened to reach the requested cap:
\nolinkurl{limited_steps} counts accepted steps that were limited before endpoint
clipping. This distinguishes the tested temporal axis from an adaptive method
that actually reused the same smaller steps at every requested resolution.

\paragraph{Implementation mapping.}
\path{core/src/mpm/grid.rs}: \nolinkurl{GridLayout::stencil}, \nolinkurl{Stencil},
\nolinkurl{Grid::deposit}. \path{core/src/mpm/particles.rs}: \nolinkurl{ParticleSet}.
\path{core/src/mpm/solver.rs}: \nolinkurl{Simulation::particle_to_grid},
\nolinkurl{grid_update}, \nolinkurl{grid_to_particle}, \nolinkurl{try_step},
\nolinkurl{stability_limit} and \nolinkurl{advance_to}.
\path{core/tests/grid_transfer.rs} tests zero-weight bookkeeping and momentum.

\section{Single-pellet coupling and energy accounting}
\label{sec:coupling}
\subsection{Impulse exchange, not a porous-mixture closure}
For a research pellet, the signed distance is the minimum distance to its
multisphere union. A grid node inside that union is constrained only when its
relative normal velocity is approaching. With node position $\mathbf{x}_i$,
\begin{equation}
 \mathbf{u}_r(\mathbf{x}_i)=\mathbf{v}_r+
       \boldsymbol\omega_r\times(\mathbf{x}_i-\mathbf{x}_r),\qquad
 \mathbf{u}_{\rm rel}=\mathbf{v}_i-\mathbf{u}_r(\mathbf{x}_i).
\end{equation}
Apply \eqref{eq:wall-projection} to $\mathbf{u}_{\rm rel}$, then add back
$\mathbf{u}_r$. If $\Delta\mathbf{p}_i=m_i(\mathbf{v}_i'-\mathbf{v}_i)$ is
the grid impulse, the pellet receives
\begin{equation}
 \Delta\mathbf{P}_r=-\sum_i\Delta\mathbf{p}_i,\qquad
 \Delta\mathbf{L}_r=-\sum_i(\mathbf{x}_i-\mathbf{x}_r)\times\Delta\mathbf{p}_i.
 \label{eq:coupling-impulse}
\end{equation}
Linear and angular action-reaction balance is explicit at this grid/pellet
exchange. The body's start-of-step velocity is used for all constrained nodes,
then the summed impulse updates the body. It is not a simultaneous finite-mass
contact solve, and the two kinetic-energy changes need not cancel or sum to a
nonpositive number. Resolution-dependent leakage, work accuracy and multiple
research pellets remain acceptance gaps.

The separate research pellet uses world angular momentum and cylinder inertia;
wall contact is spring-dashpot/Coulomb contact rather than the grid projection.
The two rigid-body implementations share concepts, not identical state formats
or all contact-history details.

\subsection{Species and linear momentum residuals}
For household species $s\in\{\text{wood},\text{waste solids},\text{water}\}$,
\begin{equation}
 r_s=\sum_{c\in\{B,D,F,R,E\}}M_{s,c}
   -M_{s,0}-M_{s,\rm added}. \label{eq:species-residual}
\end{equation}
The compartments are bed, drawer, floor, removed and evaporated, summed over
boxes where appropriate. Reclassification of intact wood into fines is not a
new source. The native check bounds each residual by
$\max(\SI{1e-12}{kg},10^{-6}M_{s,\rm expected})$.
This tests bookkeeping, not spatial distribution or transport accuracy.

For the research continuum and optional pellet, let $\mathbf{P}$ be the sum of
particle translational and pellet linear momentum. The diagnostic is
\begin{equation}
 \mathbf{r}_P=\mathbf{P}(t)+\mathbf{P}_{\rm out}(t)-\mathbf{P}(0)
 -\mathbf{I}_g-\mathbf{I}_{\rm wall}-\mathbf{I}_{r,\rm wall}.
\end{equation}
The internal coupling impulse is not subtracted a second time. Current closed
fixtures reject nonzero outflow even when it balances this ledger. Consequently,
``mass conserved including escaped material'' cannot conceal a wall leak.

\subsection{Why a stationary wall's projection loss is not external work}
For one normal wall projection, record
\begin{equation}
 \Delta E_N=\tfrac12m_i(\|\mathbf{v}^{N}\|^2-\|\mathbf{v}^*\|^2)
          =-\tfrac12m_i v_n^2\leq0,\qquad
 \Delta D_F=\tfrac12m_i(\|\mathbf{t}\|^2-\|\mathbf{v}^{+}\|^2)\geq0.
 \label{eq:wall-energy}
\end{equation}
Thus $\Delta K_{\rm grid}=\Delta E_N-\Delta D_F$ exactly, up to roundoff.
A stationary wall has zero external boundary work, since its boundary velocity
is zero. $E_N$ mixes impact-like kinetic loss with numerical operator splitting;
calling it wall work is physically misleading.

For example, a resting node given gravity velocity $-g\Delta t$ and then
projected back to zero loses $m_i g^2\Delta t^2/2$ on the transient grid each
step. Over duration $T$, this is $m_i g^2T\Delta t/2$, which vanishes linearly
with timestep. It does not mean the stationary wall physically absorbed that
amount of gravitational work from an unmoving particle.

The recorded mechanical energy is
\begin{equation}
 E_{\rm mech}=\sum_p\left[\tfrac12m_p\|\mathbf{v}_p\|^2
       -m_p\mathbf{g}\cdot\mathbf{x}_p+V_{0p}\Psi_p\right]
       +K_r-m_r\mathbf{g}\cdot\mathbf{x}_r. \label{eq:mechanical-energy}
\end{equation}
Here $K_r$ includes rigid rotation, while the particle kinetic sum does
\emph{not} include APIC affine kinetic modes. The accumulated diagnostic sum is
\begin{align}
 L_E&=E_N-D_F+W_{r,\rm contact}
      +\Delta K_{\rm coupling,grid}+\Delta K_{\rm coupling,rigid}-D_P,\nonumber\\
 r_E&=E_{\rm mech}(t)-E_{\rm mech}(0)-L_E. \label{eq:energy-residual}
\end{align}
$D_P\geq0$ is the constitutive plastic-loss estimate.
$W_{r,\rm contact}$ is the start-of-step contact-force work estimate; it includes
recoverable penalty-spring effects and is not pure dissipation. The coupling
terms are measured kinetic changes, not an invented cancellation assumption.

\textbf{Equation~\eqref{eq:energy-residual} is not a closed physical first law.}
Missing terms include particle/grid transfer and affine energy, liquid viscous
loss, complete contact spring/rotation bookkeeping and any outflow energy.
Gravity discretization can also contribute a defect. A small $r_E$ alone could
reflect cancellation; a nonzero value cannot simply be called conservation
success. Future energy acceptance needs a consistently discretized balance,
not a renaming of the residual.

\paragraph{Implementation mapping.}
\path{core/src/mpm/rigid.rs}: \nolinkurl{Pellet::signed_distance},
\nolinkurl{couple_grid}, \nolinkurl{Pellet::apply_impulse}.
\path{core/src/mpm/solver.rs}: \nolinkurl{Ledger},
\nolinkurl{Ledger::ledgered_energy}, \nolinkurl{Simulation::energy_residual},
\nolinkurl{Simulation::momentum_residual}.
\path{core/src/research.rs}: \nolinkurl{check_closed_state} and restore validation.
The restored initial mechanical-energy baseline is checked against a rebuilt
initial state; missing histories are rejected rather than fabricated as zero.
